% Source: https://artofproblemsolving.com/community/c3h1484642_astonishing_probabilities
% Stable identity: aops-topic-1484642-q01
Imagine the $10000$ people. On average $100$ have the disease and the test is positive for $97$ of them; $9900$ do not have it and the test is (wrongly) positive for $2\%$ of them, that is, $198$. Among the $97+198=295$ positive results, $97$ belong to people with the disease, so
\[
P(\text{disease}\mid\text{positive})=\dfrac{0.01\cdot0.97}{0.01\cdot0.97+0.99\cdot0.02}=\boxed{\dfrac{97}{295}}\approx32.9\%.
\]
The classmate is wrong: the population size cancels from the formula. With $N$ people the counts are $0.0097N$ and $0.0198N$, and the ratio is $\dfrac{97}{295}$ for every $N$. Only the prevalence and the accuracies matter.
